\documentclass[11pt,twoside]{book}
\usepackage[paperwidth=6in,paperheight=9in,margin=0.75in,inner=0.875in]{geometry}
\usepackage{lmodern}
\usepackage[T1]{fontenc}
\usepackage{microtype}
\usepackage{fancyhdr}
\usepackage{titlesec}
\usepackage{tabularx}
\usepackage{booktabs}
\usepackage{enumitem}

\pagestyle{fancy}
\fancyhf{}
\fancyhead[LE]{\small\itshape Prologue}
\fancyhead[RO]{\small\itshape Simply Automation}
\fancyfoot[C]{\thepage}
\renewcommand{\headrulewidth}{0.4pt}

\titleformat{\chapter}[display]
  {\normalfont\huge\bfseries\filcenter}
  {\vspace{-20pt}\normalfont\normalsize\scshape Prologue\vspace{10pt}\\\Large The Machine That Said ``Again''}
  {10pt}
  {\Huge}

\titlespacing*{\chapter}{0pt}{0pt}{40pt}

\titleformat{\section}
  {\normalfont\Large\bfseries}
  {\thesection}{1em}{}

\begin{document}

\chapter*{The Machine That Said ``Again''}
\addcontentsline{toc}{chapter}{Prologue — The Machine That Said “Again”}

Before there was a browser, there was repetition.
Someone lifted the same weight again. Someone copied the same symbol again. Someone moved the same material from one place to another again. Someone performed the same calculation again.
And eventually, someone looked at the work and asked a deceptively simple question:
\begin{quote}
\emph{What if the machine could do this part?}
\end{quote}
That question is older than computers.
It is older than software.
It is older than the word automation itself.
The answer has changed enormously over the centuries. The question has barely changed at all.

A mechanism can repeat a movement. A punched card can encode a pattern. A computer can execute instructions. A script can replay a workflow. A browser driver can click a button. A mobile automation framework can operate a phone. A pipeline can run tests while everyone sleeps. An API can bypass the interface entirely. An AI agent can now be given a goal and choose a sequence of actions.

Each step looks revolutionary when it arrives.
In retrospect, they look strangely familiar.
That is the story this book follows.
Not the complete history of automation. That would require several books, and probably a larger definition of the word automation than anyone could comfortably agree on.
Instead, this book follows seven tools and the ideas surrounding them.

Some became dominant. Some became infrastructure. Some disappeared. Some were absorbed into something else. Some remain surprisingly relevant.
They were chosen not because they are the only important tools, or even necessarily the most important tools.
They are useful because they sit at interesting moments in the story.
They expose the arguments.
They reveal the trade-offs.
And, occasionally, they become the center of battles that seem to be about software but are really about something deeper:
\begin{quote}
\textbf{Who gets to control the machine?}
\end{quote}

Consider the history of software testing.
At one point, enterprise automation meant expensive commercial platforms, record-and-playback systems, vendor support, and carefully constructed environments.
Then came open-source browser automation.
Suddenly, automation could look less like a product you purchased and more like software you wrote.
That difference mattered.
It changed who could participate. It changed how teams worked. It changed what automation engineers were expected to know. It changed the economics of the industry.

Then the browser became only one surface.
Mobile arrived.
APIs became central.
Continuous integration turned automation into infrastructure.
Modern browser tools began paying much more attention to state, waiting, isolation, tracing, and evidence.
And then something stranger happened.
The user of the automation system started changing.
For decades, we wrote instructions for machines.
Now we increasingly write tools that allow other machines to generate instructions.
That sounds like a small distinction.
It isn't.

A traditional automation script might say:
\begin{quote}
Click this button.\\
Enter this value.\\
Submit the form.\\
Verify the result.
\end{quote}
An agent might instead receive:
\begin{quote}
\emph{Complete the checkout.}
\end{quote}
The first is an instruction sequence.
The second is a goal.
Between those two statements lies one of the biggest changes in the history of automation.

But there is a temptation here.
Whenever a new automation technology appears, we tend to describe it as if history has finally arrived at its destination.
The newest tool is smarter.
The newest framework is faster.
The newest interface is simpler.
The newest agent is autonomous.

And then, usually sooner than expected, reality arrives.
The environment changes.
The application changes.
The network behaves differently.
The test passes for the wrong reason.
The locator breaks.
The state is unexpected.
The pipeline is green while something is wrong.
The agent does exactly what it was asked to do and completely misses what was meant.

Automation has always had this problem.
The instructions are precise. The world isn't.
That may be the most important sentence in this book.
Because automation history is not simply the story of machines becoming more capable.
It is the story of humans discovering, one layer at a time, that the world contains more ambiguity than their instructions accounted for.

The first solution was repetition.
Then programmability.
Then abstraction.
Then integration.
Then continuous execution.
Then observation.
Then context.
And now, increasingly, delegation.

The machine keeps moving upward.
So does the human.
The human who once operated the mechanism becomes the person who programs it.
The programmer becomes the person who designs the system.
The system designer becomes the person who sets the goals, defines the boundaries, and asks whether the machine is solving the right problem.

That is where this story begins.

\end{document}
